\documentclass[11pt,letterpaper]{article}
\usepackage[T1]{fontenc}
\usepackage{lmodern,microtype}
\usepackage[margin=1in]{geometry}
\usepackage{amsmath,amssymb,amsthm,mathtools,mathrsfs,booktabs,enumitem,xurl,xcolor}
\usepackage[colorlinks=true,linkcolor=blue!40!black,citecolor=blue!40!black,urlcolor=blue!40!black]{hyperref}
\hypersetup{pdftitle={Topological Zero Transport and the Phase of the Discriminant},pdfauthor={Drafted for Edward Baker}}
\numberwithin{equation}{section}
\newtheorem{theorem}{Theorem}[section]
\newtheorem{proposition}[theorem]{Proposition}
\newtheorem{lemma}[theorem]{Lemma}
\newtheorem{corollary}[theorem]{Corollary}
\theoremstyle{remark}\newtheorem{remark}[theorem]{Remark}
\newcommand{\C}{\mathbb C}\newcommand{\R}{\mathbb R}
\newcommand{\im}{\operatorname{Im}}\newcommand{\re}{\operatorname{Re}}
\newcommand{\HH}{\mathcal H}\newcommand{\Aconn}{\mathcal A}
\newcommand{\CC}{\mathcal C}\newcommand{\DD}{\mathscr D}
\newcommand{\UU}{\mathscr U}\newcommand{\II}{\mathcal I}
\newcommand{\JJ}{\mathcal J}\newcommand{\QQ}{\mathcal Q}
\newcommand{\RR}{\mathcal R}\newcommand{\dd}{\,d}
\newcommand{\ord}{\operatorname{ord}}\newcommand{\Log}{\operatorname{Log}}
\setlength{\emergencystretch}{2em}
\setlist{itemsep=3pt,topsep=5pt}
\allowdisplaybreaks[2]
\title{Topological Zero Transport\\and the Phase of the Discriminant}
\author{Drafted for Edward Baker}
\date{23 October 2026}
\begin{document}
\maketitle
\begin{abstract}
A multiple zero of an analytic backward heat flow carries a positive
local degree. Integrating the normalized first jet over a complete
spacetime boundary therefore relates an initial integral to later
collision charge. We evaluate the terminal integral on moving root
tubes as the phase increment of one scalar function: the squared
Vandermonde multiplied by the analytic factor evaluated at the roots.
The projective side walls have exactly zero flux. For the canonical
de Bruijn--Newman flow we prove uniform complex-time cap estimates for
the first $K=o(L\sqrt\kappa)$ roots, where $L\sim\frac14\log T$ and
$\kappa\sim8TL$. The Vandermonde phase error is linear in the population;
the analytic factor supplies the leading cap flux. Imaginary-time flux
per root is $O(1/\log T)$ on proportional late-time intervals. At a fixed
event core, shrinking the terminal caps instead concentrates half the
charge on each imaginary-time face. These estimates distinguish the
terms that become small from those that retain the topology. The
topological and algebraic ingredients are classical; the paper's focus
is their quantitative boundary realization, not a new invariant or an
RH conclusion.
\end{abstract}

\section{One boundary formula}
\label{sec:intro}

The useful question is not only where the zeros of a heat-evolved
function lie, but how their creation can be measured from a boundary.
A boundary formula separates that question into concrete analytic
tasks: estimate a terminal cap, choose side walls with manageable
flux, and account for the faces needed to close the domain. This paper
develops that separation for the actual first jet $(H,H_z)$.

Let $\tau=t+is$ and $z=x+iy$, and suppose that $H$ is holomorphic and
satisfies $H_\tau=-H_{zz}$. Away from joint zeros of $(H,H_z)$, put
\begin{equation}
 G=(H,H_z),\qquad
 \Aconn=\frac{\im(G^*dG)}{|G|^2},\qquad
 \CC=\frac{\Aconn\wedge d\Aconn}{4\pi^2}.
 \label{eq:forms}
\end{equation}
Here ``physical jet'' means that the two components have precisely this
relative normalization. Multiplying both by a positive real scalar
does not change the normalized jet. Removing a varying common phase,
or rescaling only $H_z$, generally changes individual boundary integrals.

Write $\II_t$ for a forward-oriented real-time cap integral, $\JJ_\pm$
for the outward integrals on $s=\pm Y$, and $\RR$ for the remaining
outward spatial and joining faces. The initial outward face is
$-\II_0$. When the terminal cap consists of $m$ regular root disks
$|H/H_z|\le\epsilon$, the formula to keep in view is
\begin{equation}
 \boxed{\displaystyle
 \II_0=
 \frac{k_\epsilon}{2\pi}\im\int_{T-iY}^{T+iY}
             d\log(\Delta\UU)
 +\JJ_++\JJ_-+\RR-\QQ,\qquad
 k_\epsilon=\frac{\epsilon^2}{1+\epsilon^2}.}
 \label{eq:organizing}
\end{equation}
Its scalar and integer terms are
\begin{equation}
 \Delta=\prod_{i<j}(z_j-z_i)^2,\qquad
 \UU=\prod_{j=1}^m U(\tau,z_j(\tau)),\qquad
 \QQ=\sum_{p\in W:\,H(p)=H_z(p)=0}\binom{m_p}{2},
 \label{eq:terms}
\end{equation}
where $H=U\prod_j(z-z_j)$ near the selected roots and $m_p$ is the
spatial multiplicity at an interior event. The logarithmic integral
uses a continuous phase, so it retains every winding. A complete
statement of the geometric hypotheses is Theorem~\ref{thm:transport}.

Equation~\eqref{eq:organizing} is an equality of genuine
three-dimensional boundary integrals. It does not replace the initial
cap by a one-dimensional integral on the real axis. It counts all
complex events inside the chosen four-dimensional domain, including
the real events at which new real zeros appear. At a real event of
order $m$, the number of new real zeros is $2\lfloor m/2\rfloor$;
its complex charge is $\binom m2$. These are different quantities.
If the initial slice contains a joint zero, its indentation or a local
core must be specified and retained. The theorem does not assign a
value to an unspecified singular initial cap.

The three main conclusions are as follows. The terminal cap is exactly
the phase of $\Delta\UU$, with no derivative or volume loss. A fixed
projective radius on a moving root tube makes its spatial wall flux
zero. Finally, the remaining imaginary-time faces admit quantitative
asymptotics: they are small per root on a late interval $[\theta T,T]$,
but retain collision charge when attached to a fixed event core.
The distinction matters even when only a relative bound is sought.

\subsection{The canonical quantitative result}

We use the normalization
\begin{align}
 H(\tau,z)&=\int_0^\infty e^{\tau u^2}\Phi(u)\cos(zu)\dd u,
 \label{eq:theta}\\
 \Phi(u)&=\sum_{n\ge1}
 (2\pi^2n^4e^{9u}-3\pi n^2e^{5u})e^{-\pi n^2e^{4u}}.
 \nonumber
\end{align}
Thus $H(0,z)=\xi(\frac12+\frac{iz}{2})/8$.
Define the large real saddle $L=L_T$ and its curvature by
\begin{equation}
 4\pi e^{4L}=2TL+9,\qquad
 \kappa=16\pi e^{4L}-2T=8TL-2T+36.
 \label{eq:saddle}
\end{equation}
For fixed $Y>0$, $0<r<1/4$, and an integer $K=K(T)$, set
\begin{equation}
 \varepsilon_T=\frac{K+1+L^2}{L\sqrt\kappa},\qquad
 \epsilon=\frac rL,
 \qquad K=o(L\sqrt\kappa).
 \label{eq:range}
\end{equation}
The two different symbols distinguish the analytic error scale
$\varepsilon_T$ from the physical tube radius $\epsilon$.

\begin{theorem}[Growing populations and their boundary flux]
\label{thm:headline}
For sufficiently large $T$, the first $K$ positive real zeros of
\eqref{eq:theta} continue to distinct simple roots on
$T+i[-Y,Y]$, with disjoint projective inverse disks
$|H/H_z|\le r/L$. For any consecutive block of $m\le K$ of these
roots, its forward physical cap satisfies
\begin{align}
 \frac{\II_{T,m}}m
 &=\frac{r^2YL^2}{\pi(L^2+r^2)}
       +O_Y\!\left(\frac{r^2\varepsilon_T}{L^2}\right),
 \label{eq:capheadline}\\
 \frac1m\left|\frac{k_\epsilon}{2\pi}
       \im\int_{T-iY}^{T+iY}d\log\Delta\right|
 &=O_Y\!\left(\frac{r^2\varepsilon_T}{L^2}\right).
 \label{eq:vheadline}
\end{align}
All constants are independent of $K$ and $m$. The analytic-factor
contribution is the difference of these two terms; it has leading
value $r^2Y/\pi$ per root.

Fix also $0<\theta<1$. The same population has regular root tubes over
$[\theta T,T]+i[-Y,Y]$, with the same constant radius $r/L_T$ on
each entire tube. Their projective side flux vanishes, and each
imaginary-time face satisfies
\begin{equation}
 \frac{\JJ_\pm}{m}
 =-\frac{r^2Y\log(1/\theta)}{4\pi L_T}
       +o_{Y,\theta,r}(L_T^{-1}).
 \label{eq:edgeheadline}
\end{equation}
The population is held fixed while integrating on each such interval;
it may grow between choices of $T$.
\end{theorem}

Theorem~\ref{thm:headline} concerns increasing heat time, with ranks
allowed to grow. It is not a high-height theorem at fixed heat time.
For example $K=\lfloor\sqrt\kappa\rfloor$ is permitted. The number of
included zeros then tends to infinity, the cap has a finite nonzero
limit per zero, and the late imaginary-face flux tends to zero per
zero. None of these statements asserts that the unnormalized flux
tends to zero for every growing population.

The Vandermonde is essential for identifying the pair contribution,
but it is not the whole physical cap. In this regime it has no leading
phase increment at all. The dominant common phase comes from the
canonical analytic factor. Keeping both gives a short proof and avoids
trying to estimate two large cancelling force sums separately.

\subsection{Relation to previous work and scope of novelty}

The heat deformation belongs to the classical program of de Bruijn
and Newman~\cite{deBruijn,Newman}; Newman and Wu~\cite{NewmanWu}
describe its connections with statistical physics. Ki, Kim and
Lee~\cite{KKL} proved eventual reality and simplicity at large height
for every fixed positive heat time. Polymath~\cite{Polymath},
especially Theorem~1.5 and Section~9, obtained effective high-height
approximations and zero counts. Those results use a different limiting
regime from \eqref{eq:range}. We prove the needed complex-time
estimates directly from the complete positive source.

Root dynamics and pair interaction are also established territory.
Csordas, Smith and Varga~\cite{CSV} used zero-motion equations and
inverse-square sums in their work on Lehmer pairs. Rodgers and
Tao~\cite{RT}, Section~4 and Lemma~4.18, developed localized
logarithmic Hamiltonians with exterior interaction terms in proving
$\Lambda\ge0$. Dobner~\cite{Dobner} gives a different route to
Newman's conjecture and its extension. Hall and Ho~\cite{HallHo},
Proposition~3.1, place polynomial heat-root motion in the
Calogero--Moser setting. None of the finite pair identities used here
is proposed as a new dynamics. Montgomery's pair-correlation
theorem~\cite{Montgomery} concerns a different statistical limit;
weak pair convergence alone does not supply the complex phase or
uniform small-gap information needed by a physical cap.

The topological ingredients are classical as well. Whitehead's
integral formula~\cite{Whitehead} and Chern--Simons
transgression~\cite{ChernSimons} explain the boundary form in
\eqref{eq:forms}. Rudolph~\cite{Rudolph} explicitly discusses
reducing Hopf integrals by integration along fibers; therefore the
general idea of reducing a three-form to a one-form is not new.
The discriminant and resultant identities belong to classical
elimination theory~\cite{GKZ}. Papanicolaou, Kallitsi and
Smyrlis~\cite{PKS}, Theorem~4.1, prove that multiple zeros of an
entire caloric function are isolated in complex spacetime. The local
Hermite splitting used below is consistent with that result and with
the real zero-number theory of Angenent~\cite{Angenent}, after time
reversal. Pashaev and Gurkan~\cite{PashaevGurkan} already connect
Chern--Simons vortices, holomorphic linear evolution and zero dynamics;
their gauge-field construction is distinct from the normalized
first-jet boundary considered here.

Accordingly, the contribution proposed for assessment is the explicit
physical boundary formula together with its population-uniform
asymptotic estimates and complete closure accounting. The elementary
$O(m\eta)$ Vandermonde logarithm estimate in
Lemma~\ref{lem:vandermonde} and the narrow-strip source estimate are
the quantitative steps that make the growing-population statement
possible. We do not claim priority for the invariant, the
discriminant identity, root-following coordinates, or heat-zero
creation. A targeted primary-source review found the substantial
overlaps just described; it does not establish that no equivalent
quantitative theorem exists elsewhere. This is a research working
manuscript, with originality of that narrower contribution still
subject to specialist assessment.

\section{The topological balance and the moving wall}
\label{sec:boundary}

We orient $\C^2$ by $dt\wedge ds\wedge dx\wedge dy$.
All boundaries below are piecewise smooth, with the outward
orientation. Corners have zero three-dimensional measure.

\begin{lemma}[Local heat charge]
\label{lem:charge}
Suppose $H\not\equiv0$ is holomorphic near a point $p=(\tau_*,z_*)$,
satisfies $H_\tau=-H_{zz}$, and $H(\tau_*,\cdot)$ has a zero
of order $n\ge2$ at $z_*$. Then the joint zero of $(H,H_z)$ is
isolated, and its local complex degree is $\binom n2$.
\end{lemma}
\begin{proof}
Translate $p$ to zero and write the leading spatial coefficient as
$a\ne0$. The convergent heat expansion has weighted leading term
\[
 H(\tau,z)=aP_n(\tau,z)+\text{terms of weighted degree greater than }n,
 \qquad
 P_n=\sum_{j=0}^{\lfloor n/2\rfloor}
       \frac{n!(-\tau)^jz^{n-2j}}{j!(n-2j)!},
\]
where $\tau$ has weight two and $z$ weight one.
The polynomial $P_n(1,z)$ is a scaled Hermite polynomial and has
$n$ distinct real roots. One elementary proof uses its orthogonality
with a positive Gaussian weight: fewer than $n$ distinct sign changes
would produce a polynomial of degree below $n$ whose product with
$P_n$ has one sign, contradicting orthogonality.

Put $\tau=v^2$ and $z=vw$. The quotient $v^{-n}H(v^2,vw)$
extends holomorphically to $v=0$ and equals $aP_n(1,w)$ there.
The implicit function theorem gives all $n$ roots in the cluster as
$z_j(v^2)=v(r_j+O(v))$, with distinct $r_j$.
Consequently its monic Weierstrass discriminant is
\[
 \Delta(\tau)=c\tau^{n(n-1)/2}(1+o(1)),\qquad c\ne0.
\]
It is holomorphic in $\tau$, so its zero has the indicated integer
order and is isolated. This also isolates the joint zero.

For completeness, the local intersection number of $H=0$ and
$H_z=0$ equals the vanishing order of this discriminant. Indeed,
Weierstrass preparation writes $H=UP$, with $U$ a unit and $P$
monic of degree $n$. Modulo $P$, $H_z=UP_z$; multiplication by
the unit changes the resultant by a nonvanishing factor. The
resultant of $P$ and $P_z$ is its discriminant up to sign. Equivalently,
the determinant of multiplication by $P_z$ on the free
$\C\{\tau\}$-module $\C\{\tau\}[z]/(P)$ has this vanishing
order, equal to the dimension of the local intersection algebra.
A small regular perturbation counts that many positive holomorphic
intersections, hence gives the local degree.
\end{proof}

At real coefficients the same expansion shows that, just after a real
event, all $n$ local roots are simple and real. Just before it there
are $n\bmod2$ real roots. Thus the real birth weight and complex
degree stated in Section~\ref{sec:intro} follow. Positivity of the
degree is not a sign estimate on each individual boundary face.

\begin{lemma}[Complete boundary integral]
\label{lem:degree}
Let $W$ be a compact domain on whose boundary $G$ is nonzero.
Then
\begin{equation}
 \int_{\partial W}\CC
 =\deg(G/|G|\,;\partial W)
 =\sum_{p\in W:\,G(p)=0}\binom{m_p}{2}.
 \label{eq:degree}
\end{equation}
\end{lemma}
\begin{proof}
On the unit sphere in $\C^2$, the form
$\alpha=\im(\bar\zeta_1d\zeta_1+\bar\zeta_2d\zeta_2)$ satisfies
$\alpha\wedge d\alpha=2\,d\mathrm{Vol}_{S^3}$.
Its integral is $4\pi^2$. This follows by evaluating at $(1,0)$
on $(\partial_{y_1},\partial_{x_2},\partial_{y_2})$ and using
unitary invariance. The form $\CC$ is the pullback of this normalized
volume. Excise disjoint balls around the finitely many joint zeros
and apply degree additivity, then Lemma~\ref{lem:charge}.
\end{proof}

\subsection{The exact cap integral}

A regular projective root tube over a complex-time rectangle $D$
is a holomorphic inverse
\begin{equation}
 Z:D\times\{|q|\le\epsilon\}\longrightarrow\C,
 \qquad H_z(\tau,Z)\ne0,\qquad
 \frac{H(\tau,Z)}{H_z(\tau,Z)}=q,
 \label{eq:inverse}
\end{equation}
defined on a neighborhood of the displayed compact set and injective
in each $q$ fiber. The radius is constant on this whole tube.
Differentiating the inverse equation gives $Z_q\ne0$, so its
parameterization preserves the complex spatial orientation. Its
center $z(\tau)=Z(\tau,0)$ follows a simple zero.

\begin{proposition}[Phase reduction and zero side flux]
\label{prop:phase}
For $m$ disjoint tubes of the form \eqref{eq:inverse}, with a common
radius $\epsilon$, define
\begin{equation}
 h_j(\tau)=H_z(\tau,z_j(\tau)),\qquad
 \DD(\tau)=\prod_{j=1}^m h_j(\tau).
 \label{eq:D}
\end{equation}
For an oriented heat edge $e$, the integral over its root disks is
\begin{equation}
 I(e)=\frac{k_\epsilon}{2\pi}\im\int_e d\log\DD.
 \label{eq:phase}
\end{equation}
Every wall $|q|=\epsilon$ has identically zero pullback of $\CC$.
Furthermore,
\begin{equation}
 \DD=(-1)^{m(m-1)/2}\Delta\UU.
 \label{eq:factor}
\end{equation}
The factor $U=H/\prod_j(z-z_j)$ need only be a nonvanishing
holomorphic factor near the selected roots. It can have zeros outside
those neighborhoods; no infinite product is required.
\end{proposition}
\begin{proof}
In one tube write $h(\tau,q)=H_z(\tau,Z(\tau,q))$ and
$\phi=\im\log h$, using a continuous logarithm on the contractible
parameter domain. The physical jet is exactly $h(q,1)$, so
\begin{align}
 \Aconn&=d\phi+a(q),&
 a(q)&=\frac{\im(\bar q\,dq)}{1+|q|^2},\nonumber\\
 d\Aconn&=\omega(q),&
 \omega(q)&=\frac{2\,du\wedge dv}{(1+u^2+v^2)^2},\quad q=u+iv.
 \label{eq:chart}
\end{align}
Since $a\wedge\omega=0$,
$\CC=d\phi\wedge\omega/(4\pi^2)$. The derivative of $\phi$
in the edge direction is harmonic in $q$. Its weighted average over
the disk is its center value times
\[
 \int_{|q|\le\epsilon}\omega
 =\int_0^{2\pi}\int_0^\epsilon\frac{2r}{(1+r^2)^2}\dd r\dd\vartheta
 =2\pi k_\epsilon.
\]
This proves \eqref{eq:phase} after summing the tubes.
On the wall $q=\epsilon e^{i\vartheta}$, the two-form $\omega$
pulls back to zero, proving the side assertion.

At a selected simple root,
$H_z(\tau,z_j)=U(\tau,z_j)\prod_{i\ne j}(z_j-z_i)$.
Multiplying these identities pairs the factors belonging to every
unordered pair and gives \eqref{eq:factor}.
\end{proof}

The vanishing side flux is a geometric statement, not a bound on the
speed of a root. It remains true when the disks move rapidly. A
radius allowed to vary \emph{along} the tube would produce
$dr\wedge d\vartheta$ terms and would not have the same property.
We allow a different constant radius for each chosen terminal $T$.

\begin{theorem}[The complete transport formula]
\label{thm:transport}
Suppose $W$ is an event-containing compact domain with $G\ne0$ on
its boundary, whose terminal face is the union of $m$ disks from
Proposition~\ref{prop:phase}. Decompose its other boundary into the
initial face, the two imaginary-time faces, regular projective walls,
and all remaining spatial or joining faces. Their oriented integrals
are $-\II_0$, $\JJ_+$, $\JJ_-$, $0$, and $\RR$, respectively.
Then \eqref{eq:organizing} holds with \eqref{eq:terms}.

If $A(W)$ is the number of distinct real creation events in the
interior, and separately proved bounds give
\[
 \II_0\ge L_0,\quad \II_T\le U_T,\quad
 \JJ_++\JJ_-\le U_J,\quad \RR\le U_R,
\]
then
\begin{equation}
 A(W)\le\QQ\le U_T+U_J+U_R-L_0.
 \label{eq:eventbound}
\end{equation}
\end{theorem}
\begin{proof}
Lemma~\ref{lem:degree} gives
$\QQ=\II_T-\II_0+\JJ_++\JJ_-+\RR$.
Substitute Proposition~\ref{prop:phase}. Each real event contributes
at least one to the nonnegative sum defining $\QQ$, proving the bound.
\end{proof}

This is a useful interface for estimates, but it does not prove that
the four required one-sided bounds have favorable sizes. In
particular, a bound for a closure expressed in terms of $\QQ$ itself
characterizes the closure; it cannot be inserted into
\eqref{eq:eventbound} as independent information about $\QQ$.

\subsection{The other two faces of a regular tube}

Assume real type, meaning
$H(\bar\tau,\bar z)=\overline{H(\tau,z)}$, and let the root branches be real at real time on
$[a,b]$. Each $h_j(t)$ then has a constant real sign. Choose
\begin{equation}
 \Theta(t,Y)=\sum_j\im\log\frac{h_j(t+iY)}{h_j(t)},
 \label{eq:Theta}
\end{equation}
where each logarithm is continued vertically from zero imaginary
time. Conjugation and Proposition~\ref{prop:phase} give
\begin{equation}
 \boxed{\displaystyle
 \II_t=\frac{k_\epsilon}{\pi}\Theta(t,Y),\qquad
 \JJ_+=\JJ_-=
 \frac{k_\epsilon}{2\pi}\{\Theta(a,Y)-\Theta(b,Y)\}
 =\frac{\II_a-\II_b}{2}.}
 \label{eq:closures}
\end{equation}
The orientation on the upper imaginary face is
$-dt\wedge du\wedge dv$; on the lower it is
$dt\wedge du\wedge dv$. Thus the complete regular tube obeys
$\II_b-\II_a+\JJ_++\JJ_-=0$, as it must because it contains
no joint zero. Attaching tubes to an event region does not remove
the complement of their attachment disks.

\section{A Vandermonde estimate with linear population cost}
\label{sec:vandermonde}

A termwise estimate of pair phases suggests an unnecessary harmonic
sum. Summing the linear part first removes it. This observation is
useful because a cluster has $m(m-1)/2$ pairs, whereas its physical
cap should be compared with a population of size $m$.

\begin{lemma}[Perturbing a clock discriminant]
\label{lem:vandermonde}
Let $c_j=c_1+(j-1)\pi$, $1\le j\le m$, and let
$w_j=c_j+\delta_j$ with $|\delta_j|\le\eta\le\pi/8$.
For $m=1$ both discriminants below are one. Using the logarithm
obtained by continuously moving the points from
$c_j$ to $w_j$, one has
\begin{equation}
 \left|\log\frac{\prod_{i<j}(w_j-w_i)^2}
                      {\prod_{i<j}(c_j-c_i)^2}\right|
 \le \frac{4(m-1)}\pi\eta+\frac{8m}{9}\eta^2
 \le 2m\eta.
 \label{eq:vandermondebound}
\end{equation}
\end{lemma}
\begin{proof}
Set $u_{ij}=(\delta_j-\delta_i)/(\pi(j-i))$, so $|u_{ij}|\le1/4$.
Each pair ratio stays in a disk about one avoiding zero. Hence its
logarithm is unambiguously the power-series logarithm, and
\[
 \log(\Delta(w)/\Delta(c))
 =2\sum_{i<j}u_{ij}+2\sum_{i<j}\{\Log(1+u_{ij})-u_{ij}\}.
\]
With $H_0=0$ and $H_j=\sum_{d=1}^j d^{-1}$,
\[
 \sum_{i<j}\frac{\delta_j-\delta_i}{j-i}
 =\sum_{j=1}^m\delta_j(H_{j-1}-H_{m-j}).
\]
For each $d$, the two indicator sequences in
$H_{j-1}-H_{m-j}$ differ at $2\min(d,m-d)$ indices. Therefore
\[
 \sum_{j=1}^m|H_{j-1}-H_{m-j}|
 \le\sum_{d=1}^{m-1}\frac{2\min(d,m-d)}d
 \le2(m-1).
\]
This bounds the linear part by $4(m-1)\eta/\pi$.
The elementary remainder estimate
$|\Log(1+u)-u|\le |u|^2/[2(1-|u|)]$ bounds the rest by
\[
 \frac{16\eta^2}{3\pi^2}
       \sum_{d=1}^{m-1}\frac{m-d}{d^2}
 \le\frac{8m}{9}\eta^2.
\]
Finally $4/\pi+\pi/9<2$, proving the last inequality.
\end{proof}

The estimate controls a continuously lifted complex logarithm, not
only the modulus of a determinant. A lower gap bound by itself
would not give the same phase information. The proof is elementary;
no optimality or priority for its numerical constants is asserted.

\section{The complete theta source on a narrow spatial strip}
\label{sec:source}

We now prove Theorem~\ref{thm:headline}. The central simplification
is to work in disks around the clock points, whose imaginary parts
are bounded independently of the rank. On such disks the exponential
bound for sine and cosine involves $|\im w|$, not $|w|$. The real
width can therefore grow without introducing an $e^{\pi K}$ loss.

\begin{lemma}[Weighted concentration]
\label{lem:moments}
Let $U$ be the positive random variable with density
$e^{Tu^2}\Phi(u)/H(T,0)$ on $u\ge0$, and put $V=U-L$.
For each fixed $a\ge0$ and $j=1,2,3$,
\begin{equation}
 M_j(T,a):=\mathbb E\bigl[|V|^j e^{a|V|/L}\bigr]
 =O_a(\kappa^{-j/2}).
 \label{eq:moments}
\end{equation}
The constants do not depend on a root rank or a selected population.
\end{lemma}
\begin{proof}
All theta indices can be retained in the factorization
\begin{align*}
 e^{Tu^2}\Phi(u)&=2\pi^2 e^{q_T(u)}b(u),&
 q_T(u)&=Tu^2+9u-\pi e^{4u},\\
 b(u)&=\sum_{n\ge1}\left(n^4-\frac{3n^2e^{-4u}}{2\pi}\right)
          e^{-\pi(n^2-1)e^{4u}}.
\end{align*}
For all $u\ge0$,
\begin{equation}
 0<b_*:=1-\frac3{2\pi}\le b(u)
 \le B_*:=\frac1{1-16e^{-3\pi}}.
 \label{eq:sourcebounds}
\end{equation}
The lower bound uses the first term; all terms are positive. For
the upper bound discard the subtraction and compare successive terms
$n^4e^{-\pi(n^2-1)}$ with a geometric series of ratio at most
$16e^{-3\pi}$.

For large $T$, $L\to\infty$, $T/\kappa\to0$, and
\[
 q_T(L+v)-q_T(L)=Tv^2-\pi e^{4L}(e^{4v}-1-4v).
\]
On $v\ge-1/4$ one has $q_T''(L+v)\le-\kappa/(2e)$
eventually. Hence this difference is at most
$-c\kappa v^2$, where $c=1/(4e)$.
On $0\le u\le L-1/4$, $q_T$ has at most one local minimum and
no local maximum: $q_T'$ first increases and then decreases, with
its later zero at $L$. Since $q_T(L-1/4)>q_T(0)$ eventually,
the maximum on that interval is its right endpoint. Thus the entire
far-left part has exponent at most $q_T(L)-c\kappa/16$.

On $|v|\le\kappa^{-1/2}$, Taylor's formula gives
$q_T(L+v)-q_T(L)\ge-1/2-O(\kappa^{-1/2})$.
The normalization is therefore at least a positive constant times
$e^{q_T(L)}/\sqrt\kappa$. On the near and right parts put
$w=\sqrt\kappa v$. Dividing by this lower normalization bounds
the moment by a constant times
\[
 \kappa^{-j/2}\int_\R |w|^j
 e^{-cw^2+a|w|/(L\sqrt\kappa)}\dd w
 +L^{j+1}\sqrt\kappa\,e^{a-c\kappa/16}.
\]
The integral is bounded for large $T$; the second term is smaller
than every negative power of $\kappa$. This proves the assertion.
Both tails, the normalization, and all source indices were included.
\end{proof}

\begin{lemma}[A rank-uniform complex-time approximation]
\label{lem:strip}
For real $|s|\le Y$ define
\[
 F(s,w)=\frac{H(T+is,w/L)}{H(T,0)},\qquad
 c_j(w)=\partial_w^j\cos w.
\]
On $|\im w|\le a$ and $|w|\le Z$, for $j=0,1,2,3$,
\begin{equation}
 |\partial_w^jF-e^{isL^2}c_j(w)|\le
 e^a\left\{
 \sum_{k=1}^j\binom jk\frac{M_k(T,a)}{L^k}
 +\frac{ZM_1(T,a)}L
 +Y\bigl(2LM_1(T,a)+M_2(T,a)\bigr)\right\}.
 \label{eq:stripbound}
\end{equation}
The sum is zero for $j=0$. In particular, for fixed $a,Y$ and
$Z\le\pi K+1$, the right side is $O_{a,Y}(\varepsilon_T)$.
\end{lemma}
\begin{proof}
Put $R=U/L$. Differentiation under the convergent complete source
gives $\partial_w^jF=\mathbb E[R^je^{isU^2}c_j(Rw)]$.
Split the error into the changes in $R^j$, in $c_j(Rw)$, and in
$e^{isU^2}$. The first is bounded using
\[
 |R^j-1|\le\sum_{k=1}^j\binom jk|V|^k/L^k,
 \qquad |c_j(Rw)|\le e^a e^{a|V|/L}.
\]
The line segment from $w$ to $Rw$ gives
$|c_j(Rw)-c_j(w)|\le Z|V|e^ae^{a|V|/L}/L$.
Finally, because $s,U,L$ are real,
\[
 |e^{isU^2}-e^{isL^2}|
 \le |s|\,|U^2-L^2|\le Y(2L|V|+V^2).
\]
Taking expectations proves \eqref{eq:stripbound}. No complex-time
derivative of an asymptotic error has been taken.
\end{proof}

\subsection{Actual moving disks and their rank labels}

Let $s_j=(j-1/2)\pi$, $r_1=1/3$, and $\rho=1/4$.
For $|v|\le r_1$ consider
\begin{equation}
 F(s,w)-vF_w(s,w)=0,
 \qquad w_{0,j}(v)=s_j+\arctan v.
 \label{eq:rootinverse}
\end{equation}
The principal branch of $\arctan$ is analytic here. On
$|w-w_{0,j}(v)|=\rho$, the reference equation is
\[
 e^{isL^2}(\cos w+v\sin w)
 =-\sin(s_j)e^{isL^2}\sqrt{1+v^2}\,
       \sin(w-w_{0,j}(v)).
\]
Its modulus is at least
$\sqrt{1-r_1^2}(2\rho-\sinh\rho)>0$, uniformly in $j$.
These disks have $|\im w|\le\operatorname{arctanh}(1/3)+1/4<1$
and $|w|\le\pi K+1$ for $j\le K$.
Lemma~\ref{lem:strip} and Rouch\'e's theorem therefore give exactly
one simple solution in each disk, once $\varepsilon_T$ is small.
Throughout the disk,
\[
 |\sin w|=|\cos(w-s_j)|
 \ge2-\cosh\bigl(\operatorname{arctanh}(1/3)+1/4\bigr)>0.
\]
Thus $F_w\ne0$ at the actual solutions as well.
The solutions depend holomorphically on $\tau=T+is$ and $v$
locally, and uniqueness patches them along the compact strip.
Since $H/H_z=v/L$, they are the required inverse disks in the
physical coordinate $q=v/L$.
Injectivity in $v$ follows by recovering $v$ from that quotient.
Their images are disjoint, since the disks around adjacent clock
centers have total diameter less than $\pi$.

At $s=0$ the source and the reference are real. Choose any fixed
$0<\beta<\pi/4$. On the real interval from zero to $s_K+\beta$,
the uniform errors for $F$ and $F_w$ tend to zero. On each
$[s_j-\beta,s_j+\beta]$ the derivative has constant nonzero
sign and the two endpoint values have opposite signs. On the
remaining intervals $|\cos w|\ge\sin\beta$, so no zero is
missed. This identifies the roots as exactly the first $K$
positive real roots, not merely $K$ roots near guessed labels.

At the centers $v=0$, the same estimates imply, uniformly for
$j\le K$ and $|s|\le Y$,
\begin{align}
 Lz_j(T+is)&=s_j+\delta_j(s),& |\delta_j(s)|&\le C_Y\varepsilon_T,
 \label{eq:rooterror}\\
 h_j(T+is)&=H(T,0)L[-\sin(s_j)]e^{isL^2}(1+\rho_j(s)),&
 |\rho_j(s)|&\le C_Y\varepsilon_T.
 \label{eq:herror}
\end{align}
For \eqref{eq:rooterror}, use the lower linear bound for $\sin$
on the fixed root disk. For \eqref{eq:herror}, evaluate the
$F_w$ estimate there and use the bounded derivative of $\sin$.
Increasing $C_Y$ if needed pays both bounds; it is independent of
the population. The ratio in \eqref{eq:herror} remains in a disk
about one of radius less than $1/2$, so it has a single near-one
logarithm throughout the cap. This is what controls a lifted phase,
even though $YL^2$ itself becomes large.

\subsection{Proof of the terminal estimates}

Equation~\eqref{eq:herror}, summed over the selected block, gives
\begin{equation}
 \Theta(T,Y)=mYL^2+O_Y(m\varepsilon_T).
 \label{eq:thetaerror}
\end{equation}
Substitution in \eqref{eq:closures}, with
$k_\epsilon=r^2/(L^2+r^2)$, proves \eqref{eq:capheadline}.
This argument uses endpoint logarithms and the exact disk integral;
it does not multiply an uncontrolled derivative error by a growing
cap volume.

Apply Lemma~\ref{lem:vandermonde} to the unfolded roots in
\eqref{eq:rooterror}. The reference discriminant is positive real
and independent of $s$; the factor $L^{-m(m-1)}$ is also positive
real. For $\eta=C_Y\varepsilon_T\le\pi/8$ the two endpoint
logarithms have modulus at most $2m\eta$ each. Therefore
\begin{equation}
 \left|\frac{k_\epsilon}{2\pi}
       \im\int_{T-iY}^{T+iY}d\log\Delta\right|
 \le\frac{2k_\epsilon m\eta}{\pi}.
 \label{eq:vfinite}
\end{equation}
This proves \eqref{eq:vheadline}. The factorization
\eqref{eq:factor} proves the corresponding analytic-factor estimate
by subtraction. In particular,
\begin{equation}
 \frac{k_\epsilon}{2\pi m}
       \im\int_{T-iY}^{T+iY}d\log\UU
 =\frac{r^2YL^2}{\pi(L^2+r^2)}
       +O_Y\!\left(\frac{r^2\varepsilon_T}{L^2}\right).
 \label{eq:unitbound}
\end{equation}

If $Y=Y_T=S/L^2$ with fixed $S>0$, repeat
\eqref{eq:stripbound} with that value, rather than using a fixed-width
error. Its scale becomes $O_S((K+1)/(L\sqrt\kappa))$.
The same construction gives
\begin{equation}
 \frac{\II_{T,m}}m
 =\frac{r^2S}{\pi(L^2+r^2)}
 +O_S\!\left(\frac{r^2(K+1)}{L^3\sqrt\kappa}\right)
 =O_{r,S}(L^{-2}).
 \label{eq:shrinkingcap}
\end{equation}
This is a uniform estimate per root even when $m$ grows. An absolute
vanishing cap follows, for example, if also $m=o(L^2)$.

\subsection{Proof of the late edge estimate}

On $[\theta T,T]$, $L_t/L_T\to1$ and
$\kappa_t/\kappa_T$ stays bounded above and below by positive
constants depending on $\theta$. The preceding error estimates
are therefore uniform there. For the fixed physical radius
$\epsilon=r/L_T$, the local inverse variable is
$v=L_tq$, with $|v|\le r$. The same construction gives disjoint
holomorphic tubes over the whole rectangle: the unique inverse
branches agree on overlaps. No assertion about continuation back
to time zero is needed.

At each endpoint $b\in\{\theta T,T\}$, formula
\eqref{eq:thetaerror} reads
$\II_b/m=k_\epsilon YL_b^2/\pi+O_{Y,\theta}(k_\epsilon\varepsilon_T)$.
Thus \eqref{eq:closures} yields the more precise finite asymptotic
\begin{equation}
 \frac{\JJ_\pm}{m}
 =\frac{r^2Y(L_{\theta T}^2-L_T^2)}{2\pi(L_T^2+r^2)}
 +O_{Y,\theta}\!\left(\frac{r^2\varepsilon_T}{L_T^2}\right).
 \label{eq:lateedge}
\end{equation}
Differentiating the defining saddle equation gives the exact identity
$L_t'=2L_t/\kappa_t\sim1/(4t)$, uniformly on this interval.
Hence $L_T-L_{\theta T}=\frac14\log(1/\theta)+o(1)$.
Substitution proves \eqref{eq:edgeheadline} and completes
Theorem~\ref{thm:headline}.

These tubes contain no collision. Their late edge estimate is a
proved bound on a named part of a longer boundary, not an estimate
on the earlier connection to an event-containing region. The latter
has a different asymptotic behavior, described next.

\section{What remains when a boundary encloses an event}
\label{sec:eventclosure}

An estimate on a late regular tube does not determine the flux on a
complete boundary starting at time zero. We now give a quantitative
description of this difference. It also shows why the extra terms in
\eqref{eq:organizing} cannot be suppressed by taking a thinner strip.

Let a fixed core be
\[
 W_c(Y)=[0,a]+i[-Y,Y]\;\times\;\{|z-x_*|\le R\},
 \qquad x_*\in\R,\quad 0<Y\le Y_0.
\]
Assume real type, a nonzero jet on its boundary for these $Y$, and
that the only joint zeros in the core are a fixed finite set of real
events in its interior. Denote their total complex charge by $Q$.
The initial face, the whole face $t=a$, and the spatial wall are
uniformly separated from the joint zeros. Suppose on their union
for $|s|\le Y_0$ that
\begin{equation}
 |G|\ge b>0,\qquad \|DG\|_{\rm op}\le D.
 \label{eq:corefloor}
\end{equation}
These are compact boundary quantities, not assumptions about a floor
at the events. Assume $m$ disjoint real root tubes attach along disks
of the $t=a$ face and continue to the terminal cap. The complement
of those disks remains part of the boundary.

For the canonical flow, a fixed finite set of real simple branches
that remains simple for $t\ge a$ admits such tubes for every finite
$T$. A compact initial segment has a common positive complex-time
width and inverse radius by the implicit function theorem. The
late part is covered by Section~\ref{sec:source}. Combining them,
and decreasing the common width if necessary, permits a fixed small
$Y>0$ and radius $r/L_T$ for all sufficiently large $T$. For
example, branches starting at $a>1/2$ stay simple: de Bruijn's
reality theorem and the local splitting lemma rule out a later
collision. This continuation statement neither asserts an actual
earlier theta event nor supplies a height-uniform core floor.

\begin{theorem}[Quantitative closure allocation]
\label{thm:core}
In this fixed-core geometry let $\mathcal R_T^{\rm all}$ be the
signed sum of the initial, joining and local spatial faces, including
the initial sign $-\II_0$. Set
\begin{equation}
 B=\frac{D^3}{2\pi^2b^3}(4\pi R^2+4\pi Ra).
 \label{eq:Bcore}
\end{equation}
Every projective tube wall has zero flux, and
\begin{equation}
 |\mathcal R_T^{\rm all}|\le BY,\qquad
 \boxed{\displaystyle
 \JJ_+=\JJ_-=
 \frac{Q-\II_{T,m}-\mathcal R_T^{\rm all}}2.}
 \label{eq:coreidentity}
\end{equation}
For the following asymptotic assertions, assume that $H$ is the
canonical flow \eqref{eq:theta} and that the attached roots have
fixed positive ranks. If $Y_T=S/L_T^2$ and $\epsilon_T=r/L_T$, with
fixed $S,r$ and fixed attached population, then
\begin{equation}
 \II_{T,m}=O(mL_T^{-2}),\qquad
 \left|\JJ_\pm-\frac Q2\right|
 \le\frac{BY_T+|\II_{T,m}|}{2}
 =O((B+m)L_T^{-2}).
 \label{eq:halfcharge}
\end{equation}
For fixed admissible $Y>0$, let $R_c(Y)$ be the signed sum of the
three non-imaginary faces of the unmodified core. Then
\begin{align}
 \II_{T,m}&\longrightarrow \frac{mr^2Y}{\pi},\nonumber\\
 \mathcal R_T^{\rm all}&\longrightarrow R_c(Y),
       &|R_c(Y)|&\le BY,\nonumber\\
 \JJ_\pm&\longrightarrow
      \frac{Q-R_c(Y)}2-\frac{mr^2Y}{2\pi}.
 \label{eq:fixedwidthcore}
\end{align}
The errors in these fixed-width limits are
$O(L_T^{-2})+O(\varepsilon_T/L_T^2)$ with fixed core and
population constants.
\end{theorem}
\begin{proof}
The normalized sphere-volume expression for \eqref{eq:forms}
gives a comass bound $D^3/(2\pi^2b^3)$ on the faces in
\eqref{eq:corefloor}. Their total three-volume is at most
$Y(4\pi R^2+4\pi Ra)$. Removing the attachment disks only
reduces this upper bound. This proves the first inequality.

Simultaneous conjugation $(t,s,x,y)\mapsto(t,-s,x,-y)$ preserves
the ambient orientation, exchanges the outward imaginary faces,
and sends $\Aconn$ and $d\Aconn$ to their negatives. It therefore
preserves $\CC$ and proves $\JJ_+=\JJ_-$. The complete boundary
has charge $Q$ by Lemma~\ref{lem:degree}, yielding
\eqref{eq:coreidentity}. Equation~\eqref{eq:shrinkingcap} proves
\eqref{eq:halfcharge}.

For fixed width, the removed attachment caps at $t=a$ have size
$O(k_{\epsilon_T})=O(L_T^{-2})$, because their center derivatives
are nonvanishing on a fixed compact segment and their lifted phase
increments are bounded. Thus
$\mathcal R_T^{\rm all}=R_c(Y)+O(L_T^{-2})$.
The terminal estimate \eqref{eq:capheadline} and the exact balance
give \eqref{eq:fixedwidthcore} with the stated errors.
\end{proof}

The half-charge limit is a characterization of the remaining flux,
not an independent upper bound on the event count. Its constants
belong to a fixed core. Allowing height, multiplicity or the number
of cores to grow requires new estimates for $b,D,R,a$; none is
implicit here. Shrinking $Y$ also changes the initial cap being
integrated. At fixed $Y$, the initial core and its initial integral
remain fixed, and the nonzero limits in \eqref{eq:fixedwidthcore}
give the appropriate comparison.

\subsection{Two exact controls}

The quadratic heat solution
$H(\tau,z)=z^2-2\tau$ has charge one at the origin. Its inverse
root disks are
\[
 Z_\pm(\tau,q)=q\pm\sqrt{q^2+2\tau}.
\]
For $t\ge a>0$ and $\epsilon^2<2a$, these are disjoint regular
tubes on a sufficiently small imaginary-time strip. Their two-root
terminal cap is exactly
\begin{equation}
 \II_T=\frac{k_\epsilon}{\pi}\arctan(Y/T).
 \label{eq:quadraticcap}
\end{equation}
Indeed $\DD=-8\tau$, so the vertical phase increment is
$2\arctan(Y/T)$. The cap tends to zero and the sides have zero
flux, although an attached core still has charge one. The two
inverse branches pinch at $q^2=-2\tau$ when continued backward;
one cannot extend a fixed-radius regular disk bundle through the
event. This example checks the need for a joining region.

For $\lambda>0$, the exact heat solution
$H(\tau,z)=e^{\lambda^2\tau}\sin(\lambda z)/\lambda$
has stationary simple roots. On any fixed finite collection of
their projective disks, $\Delta$ is constant but
$\DD$ has phase rate $m\lambda^2$. Hence
\begin{equation}
 \II_T=\frac{k_\epsilon m\lambda^2Y}{\pi}.
 \label{eq:sinecap}
\end{equation}
This cap is nonzero; the full regular tube still has zero charge.
It checks that root positions alone do not determine the physical
cap. The leading phase in Theorem~\ref{thm:headline} is the
canonical, increasing-frequency counterpart of this calculation.

Both controls satisfy the heat equation but lack the unchanged
theta source. Conversely, the source proof uses only the displayed
positive density bounds and concentration: it is not claimed to
distinguish theta from every other source with the same asymptotic
properties. No collision exclusion can be inferred merely from
the late clock pattern.

\section{What the formula now makes accessible}
\label{sec:outlook}

The initial-to-future identity \eqref{eq:organizing} and its
independent-bound consumer \eqref{eq:eventbound} are the organizing
framework. The estimates proved here supply a terminal cap, its
Vandermonde and analytic-factor decomposition, exactly vanishing
projective sides, and relative late-edge asymptotics for a growing
population. They also identify a substantial surviving term at a
fixed event core. These are concrete bounds with different domains
of validity; their domains must match before they can be combined
to bound events.

The next estimate needed for a full time-zero application is on the
complete earlier attachment and imaginary-time boundary, with
constants controlled as the spatial height grows. A meaningful
first target is a bound for that contribution divided by the
enclosed zero population. It need not be small enough to exclude
every event. One could then use an independently obtained lower
bound on the fixed initial cap in \eqref{eq:eventbound}. Neither
that lower bound nor a uniform early-boundary estimate is proved
here. The fixed-core half-charge formula cannot replace them,
because it already contains the unknown charge.

There is also a separate real-variable route. At fixed positive
heat time, the high-height estimates of Polymath yield zero-free
moving barriers and finite budgets for later real births. Those
budgets do not estimate the physical imaginary-time integrals in
this paper. Likewise a pair-correlation law for real zeros does
not supply the analytic factor in \eqref{eq:organizing}. Keeping
these distinctions explicit leaves a modest but useful conclusion:
the topological formula has quantitatively tractable pieces, and
the remaining task is located on a specified boundary rather than
hidden in a presumed vanishing cap.

\section*{Acknowledgements and research status}
This manuscript was drafted for Edward Baker with substantial
assistance from OpenAI's Codex, recorded model \texttt{gpt-6-astra},
reasoning effort \emph{max}. LLM assistance was used for derivation,
literature comparison, checking and exposition. The checks in this
round are by the same assistant; they are not independent specialist
or human review. The general transport and fixed-population cap
arguments consolidate earlier LLM-assisted research notes in the
author's archive~\cite{Archive}. The present continuation adds the
narrow-strip growing-population estimates, the linear-population
Vandermonde logarithm bound and the proportional-interval edge
asymptotic. The companion working manuscript on zero counts and
localized Hodge energy remains a separate catalogue.

All asymptotic statements have analytic proofs. No numerical onset,
theta-zero certificate, priority certification, proof of RH, or
proof of initial simplicity is asserted. The accompanying literature
and proof review records source access, overlap, control examples
and the remaining novelty question.

\begin{thebibliography}{99}
\bibitem{deBruijn}
N.~G.~de Bruijn, \emph{The roots of trigonometric integrals},
Duke Mathematical Journal \textbf{17} (1950), 197--226.
\url{https://doi.org/10.1215/S0012-7094-50-01720-0}.

\bibitem{Newman}
C.~M.~Newman, \emph{Fourier transforms with only real zeros},
Proceedings of the American Mathematical Society \textbf{61}
(1976), 245--251.
\url{https://doi.org/10.1090/S0002-9939-1976-0434982-5}.

\bibitem{NewmanWu}
C.~M.~Newman and W.~Wu, \emph{Constants of de Bruijn--Newman type
in analytic number theory and statistical physics},
Bulletin of the American Mathematical Society \textbf{57}
(2020), 595--614.
\url{https://doi.org/10.1090/bull/1668}.

\bibitem{KKL}
H.~Ki, Y.-O.~Kim and J.~Lee, \emph{On the de Bruijn--Newman
constant}, Advances in Mathematics \textbf{222} (2009), 281--306.
\url{https://doi.org/10.1016/j.aim.2009.04.003}.

\bibitem{Polymath}
D.~H.~J.~Polymath, \emph{Effective approximation of heat flow
evolution of the Riemann $\xi$ function, and a new upper bound
for the de Bruijn--Newman constant}, 2019.
\url{https://arxiv.org/abs/1904.12438v2}.

\bibitem{CSV}
G.~Csordas, W.~Smith and R.~S.~Varga, \emph{Lehmer pairs of
zeros, the de Bruijn--Newman constant $\Lambda$, and the Riemann
hypothesis}, Constructive Approximation \textbf{10} (1994), 107--129.
\url{https://doi.org/10.1007/BF01205170}.

\bibitem{RT}
B.~Rodgers and T.~Tao, \emph{The de Bruijn--Newman constant is
non-negative}, Forum of Mathematics, Pi \textbf{8} (2020), e6.
Corrected version arXiv:1801.05914v5, 2021.
\url{https://arxiv.org/abs/1801.05914v5}.

\bibitem{Dobner}
A.~Dobner, \emph{A proof of Newman's conjecture for the extended
Selberg class}, Acta Arithmetica \textbf{201} (2021), 29--62.
\url{https://doi.org/10.4064/aa200603-23-7}.
Author's updated text: \url{https://arxiv.org/abs/2005.05142v2}.

\bibitem{HallHo}
B.~C.~Hall and C.-W.~Ho, \emph{The heat flow conjecture for
polynomials and random matrices}, Letters in Mathematical Physics
\textbf{115} (2025), article 60.
\url{https://doi.org/10.1007/s11005-025-01946-9}.

\bibitem{Montgomery}
H.~L.~Montgomery, \emph{The pair correlation of zeros of the zeta
function}, in \emph{Analytic Number Theory}, Proceedings of
Symposia in Pure Mathematics, vol.~24, American Mathematical
Society, 1973, 181--193.
\url{https://www-personal.umich.edu/~hlm/paircor1.pdf}.

\bibitem{Whitehead}
J.~H.~C.~Whitehead, \emph{An expression of Hopf's invariant as an
integral}, Proceedings of the National Academy of Sciences USA
\textbf{33} (1947), 117--123.
\url{https://doi.org/10.1073/pnas.33.5.117}.

\bibitem{ChernSimons}
S.-S.~Chern and J.~Simons, \emph{Characteristic forms and
geometric invariants}, Annals of Mathematics \textbf{99}
(1974), 48--69.
\url{https://doi.org/10.2307/1971013}.

\bibitem{Rudolph}
L.~Rudolph, \emph{Whitehead's integral formula, isolated critical
points, and the enhancement of the Milnor number},
Pure and Applied Mathematics Quarterly \textbf{6} (2010), 545--554.
\url{https://arxiv.org/abs/0912.4974}.

\bibitem{GKZ}
I.~M.~Gelfand, M.~M.~Kapranov and A.~V.~Zelevinsky,
\emph{Discriminants, Resultants, and Multidimensional Determinants},
Birkh\"auser, Boston, 1994.
\url{https://doi.org/10.1007/978-0-8176-4771-1}.

\bibitem{PKS}
V.~G.~Papanicolaou, E.~Kallitsi and G.~Smyrlis,
\emph{Entire solutions for the heat equation}, Electronic Journal
of Differential Equations \textbf{2021}, no.~44, 1--25.
\url{https://ejde.math.txstate.edu/Volumes/2021/44/papanicolaou.pdf}.

\bibitem{Angenent}
S.~Angenent, \emph{The zero set of a solution of a parabolic
equation}, Journal f\"ur die reine und angewandte Mathematik
\textbf{390} (1988), 79--96.
\url{https://doi.org/10.1515/crll.1988.390.79}.

\bibitem{PashaevGurkan}
O.~K.~Pashaev and Z.~N.~Gurkan, \emph{Abelian Chern--Simons
vortices and holomorphic Burgers hierarchy}, Theoretical and
Mathematical Physics \textbf{152} (2007), 1017--1029.
\url{https://doi.org/10.1007/s11232-007-0086-0}.
\url{https://arxiv.org/abs/hep-th/0611029}.

\bibitem{Archive}
E.~Baker, \emph{Shifted-zeta-positivity research archive},
working notes prepared with substantial LLM assistance, October 2026.
The source crosswalk is retained with this manuscript; pertinent
groups are Investigations~4/11 (finite root tubes and complete
closures, Notes~01--04) and Investigations~4 (initial cap transport,
Note~11; growing complex caps, Note~12).
\url{https://github.com/ebbaker/shifted-zeta-positivity}.
\end{thebibliography}

\end{document}
